\begin{ex}
Số 84 520 gồm mấy chục nghìn?
\shortans{8}
\loigiai{
Chữ số hàng chục nghìn là 8.
}
\end{ex}

\begin{ex}
Số liền sau của 89 999 là số nào?
\shortans{90000}
\loigiai{
$89999 + 1 = 90000$.
}
\end{ex}

\begin{ex}
Số liền trước của 100 000 là số nào?
\shortans{99999}
\loigiai{
$100000 - 1 = 99999$.
}
\end{ex}

\begin{ex}
Trong các số $75400, 74500, 75040, 74050$, số lớn nhất là số nào?
\shortans{75400}
\loigiai{
75400 là số lớn nhất.
}
\end{ex}

\begin{ex}
Làm tròn số 58 400 đến hàng nghìn được số nào?
\shortans{58000}
\loigiai{
Hàng trăm là 4 nên làm tròn thành 58 000.
}
\end{ex}

\begin{ex}
Làm tròn số 65 000 đến hàng chục nghìn được số nào?
\shortans{70000}
\loigiai{
Hàng nghìn là 5 nên làm tròn lên 70 000.
}
\end{ex}

\begin{ex}
Viết số thích hợp: $70000 + 5000 + 400 + 3 = \dots$
\shortans{75403}
\loigiai{
Số đó là 75403.
}
\end{ex}

\begin{ex}
Tìm $x$ là số tròn chục nghìn thỏa mãn $70000 < x < 90000$.
\shortans{80000}
\loigiai{
Số đó là 80000.
}
\end{ex}

\begin{ex}
Cho 5 chữ số 1, 0, 4, 6, 8. Số lớn nhất có 5 chữ số khác nhau ghép từ các chữ số trên là số nào?
\shortans{86410}
\loigiai{
Sắp xếp giảm dần: 86410.
}
\end{ex}

\begin{ex}
Cho 5 chữ số 2, 0, 3, 5, 7. Số bé nhất có 5 chữ số khác nhau ghép từ các chữ số trên là số nào?
\shortans{20357}
\loigiai{
Số đó là 20357.
}
\end{ex}

\begin{ex}
Khoảng cách giữa hai số 80 000 và 76 000 là bao nhiêu?
\shortans{4000}
\loigiai{
$80000 - 76000 = 4000$.
}
\end{ex}

\begin{ex}
Số 65 000 hơn số 61 500 bao nhiêu đơn vị?
\shortans{3500}
\loigiai{
$65000 - 61500 = 3500$.
}
\end{ex}

\begin{ex}
Số tự nhiên nhỏ nhất có 5 chữ số khác nhau là số nào?
\shortans{10234}
\loigiai{
Số đó là 10234.
}
\end{ex}

\begin{ex}
Số tự nhiên lớn nhất có 5 chữ số khác nhau là số nào?
\shortans{98765}
\loigiai{
Số đó là 98765.
}
\end{ex}

\begin{ex}
Làm tròn số 47 620 đến hàng nghìn ta được số nào?
\shortans{48000}
\loigiai{
Hàng trăm là 6 nên làm tròn lên 48 000.
}
\end{ex}

\begin{ex}
Làm tròn số 83 200 đến hàng chục nghìn ta được số nào?
\shortans{80000}
\loigiai{
Hàng nghìn là 3 nên làm tròn thành 80 000.
}
\end{ex}

\begin{ex}
Có bao nhiêu số tròn chục nghìn có 5 chữ số?
\shortans{9}
\loigiai{
10000, 20000, ..., 90000 (9 số).
}
\end{ex}

\begin{ex}
Điền số thích hợp: $49999 < \dots < 50001$.
\shortans{50000}
\loigiai{
Số ở giữa là 50000.
}
\end{ex}

\begin{ex}
Số 77 777 có bao nhiêu chữ số 7?
\shortans{5}
\loigiai{
Có 5 chữ số 7.
}
\end{ex}

\begin{ex}
Tính: $50000 + 40000 = \dots$
\shortans{90000}
\loigiai{
$50000 + 40000 = 90000$.
}
\end{ex}

\begin{ex}
Tính: $100000 - 30000 = \dots$
\shortans{70000}
\loigiai{
$100000 - 30000 = 70000$.
}
\end{ex}

\begin{ex}
Tính: $20000 \times 4 = \dots$
\shortans{80000}
\loigiai{
$20000 \times 4 = 80000$.
}
\end{ex}

\begin{ex}
Tính: $90000 : 3 = \dots$
\shortans{30000}
\loigiai{
$90000 : 3 = 30000$.
}
\end{ex}

\begin{ex}
Làm tròn số 94 800 đến hàng chục nghìn được số nào?
\shortans{90000}
\loigiai{
Hàng nghìn là 4 nên làm tròn thành 90 000.
}
\end{ex}

\begin{ex}
Làm tròn số 95 000 đến hàng chục nghìn được số nào?
\shortans{100000}
\loigiai{
Hàng nghìn là 5 nên làm tròn lên 100 000.
}
\end{ex}
